\documentclass[12pt,a4paper]{article}

\usepackage[utf8]{inputenc}
\usepackage[T1]{fontenc}
\usepackage[french]{babel}
\usepackage{lmodern}
\usepackage{amsmath,amssymb}
\usepackage{array}
\usepackage{booktabs}
\usepackage{longtable}
\usepackage{geometry}
\usepackage{microtype}
\usepackage{hyperref}
\usepackage{enumitem}
\geometry{top=2.5cm,bottom=2.5cm,left=2.5cm,right=2.5cm}
\hypersetup{hidelinks,pdftitle={Documentation mathematique globale HCANN}}

\title{Documentation Globale de l'Implémentation HCANN}
\author{Projet HCANN-Agent}
\date{\today}

\begin{document}

\maketitle

\begin{abstract}
Ce document est une \textbf{vue d'ensemble} de l'implémentation HCANN.  
Il évite volontairement les redondances mathématiques et renvoie les détails complets vers les trois documents autonomes: \texttt{cortex.tex}, \texttt{hippocampus.tex}, \texttt{hcann.tex}.
\end{abstract}

\tableofcontents
\newpage

\section{Contexte général}
L'architecture HCANN (\textit{Hippocampo-Cortical Artificial Neural Network}) est conçue comme une mémoire multimodale bio-inspirée.  
Le découpage en blocs suit une logique de séparation des fonctions:
\begin{itemize}[leftmargin=*]
    \item \textbf{Cortex}: représentation sémantique et contexte court terme;
    \item \textbf{Hippocampe}: encodage rapide, dynamique de scaffold et attracteurs;
    \item \textbf{HCANN global}: orchestration rapide/lente et projections hétéro-associatives.
\end{itemize}

\section{Convention de lecture}
\begin{itemize}[leftmargin=*]
    \item $B$: taille du batch.
    \item $\mathbf{S}$: représentation sémantique corticale.
    \item $\mathbf{H}$: représentation hippocampique rapide.
    \item \textbf{Principe}: ce fichier décrit le flux global; les dérivations détaillées restent dans les trois fichiers spécialisés.
\end{itemize}

\section{Cartographie minimale des fonctions}
{\small
\setlength{\tabcolsep}{4pt}
\renewcommand{\arraystretch}{1.15}
\begin{longtable}{>{\raggedright\arraybackslash}p{2.4cm} >{\raggedright\arraybackslash}p{4.2cm} >{\raggedright\arraybackslash}p{3.0cm} >{\raggedright\arraybackslash}p{4.2cm}}
\toprule
\textbf{Fichier} & \textbf{Fonction} & \textbf{Entrées} & \textbf{Sorties} \\
\midrule
\endhead
\texttt{cortex.py} &
\texttt{MultimodalEncoder.forward(images,texts)} &
images $(B,C,H,W)$,\\ texts (batch texte) &
$\mathbf{S}$ (embedding sémantique) \\
\texttt{cortex.py} &
\texttt{WorkingMemory.forward(x)} &
$\mathbf{x}$ &
\texttt{read}, poids des slots \\
\texttt{cortex.py} &
\texttt{WorkingMemory.diversity\_loss()} &
slots internes &
scalaire $\mathcal{L}_{div}$ \\
\texttt{cortex.py} &
\texttt{EntorhinalGateway.forward(attractor,wm)} &
état attracteur + état WM &
projection vers espace hippocampique \\
\texttt{hippocampus.py} &
\texttt{GridCellModule.forward(velocity)} &
$\mathbf{u}$ &
code grille \\
\texttt{hippocampus.py} &
\texttt{HippocampalScaffold.forward(velocity,steps)} &
vitesse + nombre d'itérations &
$\mathbf{H}$ et état grille concaténé \\
\texttt{hippocampus.py} &
\texttt{HippocampalScaffold.reset\_phases()} &
aucune &
phases remises à zéro \\
\texttt{hcann.py} &
\texttt{HCANN.fast\_encode(images,texts,velocity)} &
modalités + vitesse &
$\mathbf{H}_{fast}$, $\mathbf{S}$, $\mathbf{M}$ \\
\texttt{hcann.py} &
\texttt{HCANN.slow\_consolidate(hpc\_target,sem\_input)} &
$\mathbf{H}_{target}, \mathbf{S}_{in}$ &
perte $\mathcal{L}_{hetero}$ \\
\bottomrule
\end{longtable}
}

\section{Synthèse inter-blocs}
Le flux d'information principal est:
\begin{enumerate}[leftmargin=*]
    \item Encodage multimodal cortical $\rightarrow$ embedding sémantique;
    \item Projection vers l'espace hippocampique + état scaffold;
    \item Fusion rapide pour la trace épisodique;
    \item Consolidation lente par reconstruction hippocampe$\rightarrow$cortex.
\end{enumerate}

Cette organisation permet de préserver la cohérence entre inspiration neuroscientifique (CLS, grille, attracteurs) et contraintes d'ingénierie (modularité, testabilité, exécution incrémentale).

\section{Équation pivot (sans redondance)}
Pour garder une vue compacte, on retient uniquement l'équation centrale de couplage:
\[
\mathbf{H}_{fast}=\mathbf{H}_{grid}+W_{s\to h}\mathbf{S}.
\]
Les développements complets (avec étapes intermédiaires, dimensions et correspondance code) sont documentés dans:
\begin{itemize}[leftmargin=*]
    \item \texttt{cortex.tex} (encodage, mémoire de travail, passerelle);
    \item \texttt{hippocampus.tex} (grid cells, attracteurs, scaffold);
    \item \texttt{hcann.tex} (orchestration fast/slow, perte de consolidation).
\end{itemize}

\section{Guide de compilation}
\begin{itemize}[leftmargin=*]
    \item \textbf{Mode autonome recommandé}: compiler séparément
    \texttt{cortex.tex}, \texttt{hippocampus.tex}, \texttt{hcann.tex}.
    \item \textbf{Raison}: chaque fichier contient déjà son propre \texttt{\textbackslash documentclass} et tous les packages nécessaires.
    \item \textbf{Rôle de ce fichier}: fournir une vue intégrée, sans répéter les preuves détaillées.
\end{itemize}

\section*{Note d'usage}
Ce document est compilable tel quel et sert de page de navigation globale.  
Les documents \texttt{cortex.tex}, \texttt{hippocampus.tex} et \texttt{hcann.tex} constituent les références détaillées pour la démonstration mathématique.

\end{document}
